\subsection{Taylor's formula}

2.5.1. Let \(r\) be an integer \(\geqslant 1\) and let \(f\) be a map of class \(C^r\) from an open set U of E into F. For \(x \in U\) , \(h \in E\) and \(p \leqslant r\) , let us agree to write \(D^{p}f(x_0).h^p\) instead of \(D^{p}f(x_0).(h,\ldots,h)\) . If the segment \([x,x + h]\) is contained in U, one has the formula ("Taylor's formula"):

\[
f (x + h) = \sum_ {p = 0} ^ {r - 1} \frac {1}{p !} \mathbf {D} ^ {p} f (x). h ^ {p} + v _ {r} (x; h)
\]

where the "remainder" \(v_{r}(x; h)\) is given by:

\[
v _ {r} (x; h) = \int_ {0} ^ {1} \frac {(1 - t) ^ {r - 1}}{(r - 1) !} \mathrm{D} ^ {r} f (x + t h). h ^ {r} d t
\]

One has:

\[
v _ {r} (x; h) \equiv \frac {1}{r !} \mathrm{D} ^ {r} f (x). h ^ {r} \quad \text {   mod   } o (\| h \| ^ {r}) \quad \text {   when   } h \text {   tends   to   } 0
\]

and

\[
f (x + h) \equiv \sum_ {p = 0} ^ {r} \frac {1}{p !} \mathrm{D} ^ {p} f (x). h ^ {p} \quad \mod o (\| h \| ^ {r})
\]

as h tends to zero.

2.5.2. Let moreover \(\gamma\) be a continuous seminorm on \(F\) ; if one has \(\|D^{r}f(z)\|_{\gamma}\leqslant M\) for every point \(z\) of the segment \([x,x + h]\) , then one has:

\[
\| v _ {r} (x; h) \| _ {\gamma} \leqslant \frac {\mathbf {M}}{r !} \| h \| ^ {r}
\]

2.5.3. Suppose in addition that \(\mathbf{E} = \mathbf{R}^n\) . One then has:

\[
f (x + h) \equiv \sum_ {| \alpha | \leqslant r} \Delta^ {\alpha} f (x) h ^ {\alpha} \quad \mod o (\| h \| ^ {r}) \text {   when   } h \text {   tends   to   } 0
\]

by setting:

\[
\Delta^ {\alpha} f (x) = \frac {1}{\alpha !} \partial^ {\alpha} f (x)
\]

2.5.4. Let \(f\) and \(g\) two functions of class \(C^r\) on an open set \(U\) of \(E\) , with values in \(F\) . For \(f\) and \(g\) to have at a point \(x\) of \(U\) a contact of order \(\geqslant r\) , it is necessary and sufficient that one has \(D^p f(x) = D^p g(x)\) for every integer \(p\) with \(0 \leqslant p \leqslant r\) . When \(E\) is of finite dimension, this amounts to saying that the iterated partial derivatives of order \(\leqslant r\) of \(f\) and of \(g\) (with respect to a basis of \(E\) ) are equal at the point \(x\) .

2.5.5. Let U be an open subset of \(E \times R^{n}\) of the form \(V \times I_{1} \times \cdots \times I_{n}\) , where V is an open subset of E and \(I_{1}, \ldots, I_{n}\) open intervals of R containing 0. Set \(U_{0} = V\) and \(U_{j} = V \times I_{1} \times \cdots \times I_{j}\) for \(1 \leqslant j \leqslant n\) . Given a function \(f \in \mathcal{C}^{r}(U; F)\) (with \(1 \leqslant r \leqslant \infty\) ), there exists one and only one sequence of functions \(f_{j}\in\mathcal{C}^{r-1}(\mathbf{U}_{j};\mathbf{F})\) (for \(0\leqslant j\leqslant n\) ) such that:

\[
f (x, t _ {1}, \dots , t _ {n}) = f _ {0} (x) + \sum_ {j = 1} ^ {n} t _ {j} f _ {j} (x, t _ {1}, \dots , t _ {j})
\]

for \(x \in V\) and \(t_j \in I_j\) . We have:

\[
f _ {0} (x) = f (x, 0, \dots , 0)
\]

\[
f _ {j} (x, t _ {1}, \dots , t _ {j}) = \int_ {0} ^ {1} \partial_ {j} f (x, t _ {1}, \dots , t _ {j - 1}, t _ {j} u, 0, \dots , 0) d u
\]

for \(1 \leqslant j \leqslant n\) . In this last formula, \(\partial_j f\) denotes the \(j\) -th partial derivative of the function \((t_1, \ldots, t_n) \mapsto f(x, t_1, \ldots, t_n)\) .

